\documentclass[10pt]{article}

% Specify the margins and text size.
\setlength{\textwidth}{6.6in}
\setlength{\textheight}{9.5in}
\setlength{\oddsidemargin}{0pt}
\setlength{\evensidemargin}{0pt}
\setlength{\topmargin}{0pt}
\setlength{\hoffset}{-.15in}
\setlength{\voffset}{-1in}

\setlength{\parskip}{5pt}
\setlength{\parindent}{0pt}

% Load some fonts and symbol packages
\usepackage{latexsym}
\usepackage{pifont}       % contains 'star' symbol for counterinsurgency handbook title
\usepackage{yfonts} 
\usepackage{amsmath}
\usepackage{amsfonts}

\usepackage{graphicx}     % actually, this is loaded by pstricks
\usepackage[T1]{fontenc}
\usepackage{ifthen}
\usepackage{pstricks,pst-grad,pst-text,pst-node,multido,pst-plot,calc,pst-func}
%\usepackage[all]{xy}
%\usepackage{animate}

% The hyperref package inserts links into the pdf.
\definecolor{MyLinkColor}{rgb}{.1,.2,1}
\definecolor{MyCiteColor}{rgb}{.1,1,.2}
\definecolor{MyURLColor}{rgb}{.4,.4,.4}
\usepackage[backref=true,pagebackref=false,hyperindex,colorlinks=true,
  linkcolor=MyLinkColor,urlcolor=MyURLColor]{hyperref}


% The tweaklist package is something I found on the web.  It provides a simple interface
% for making changes to spacing used in the itemize and enumerate environments.  Comment
% this out if you don't care to use tweaklist.
\usepackage{tweaklist}
\renewcommand{\itemhook}{\setlength{\parskip}{2pt}\setlength{\parsep}%
{1pt}\setlength{\topsep}{0pt}\setlength{\itemsep}{0pt}}

\newcommand{\U}{\underline{\hspace{5pt}}}

\usepackage{listings}

\begin{document}
\pagestyle{empty}
\lstset{language=R, showspaces=false, showstringspaces=false}

\href{http://www.su.edu}{\includegraphics[height=1.75cm]{sulogo.eps}}
\vspace{-1.69cm}

{\small{\ }\hfill
\begin{tabular}{cl}
& Math 207\\
& Introduction to Statistics\\
\end{tabular}
}
\setlength{\baselineskip}{1.05\baselineskip}

\begin{center}
\textbf{\large  Regression and the Normal Curve}
\end{center}
\newcommand{\Z}{\hphantom{0}}
\newcommand{\VP}{\vphantom{\Large Y}}
\newcommand{\M}{\hphantom{$-$}}

Suppose we have a list of points $(x_1, y_1)$, $(x_2, y_2)$, $\dots$, $(x_n, y_n)$
and we draw a line (such as the regression line or perhaps some other line) to 
try to predict $y$ from $x$.  
For each of the $x$-values
$x_i$ in our data list, we can use the  line to get a  predicted value 
 $y_{\mbox{\footnotesize predicted}}$ of $y$.  
We also have the true value $y_i$ of $y$.  
The difference\vspace{-3pt}
\[y_{\mbox{\footnotesize predicted}} \; -\; y_i\vspace{-3pt}\]
is the \textit{residual error} at $x_i$.  
The  RMS (root-mean-squared) size of these errors is:\vspace{-3pt}
\[\mbox{RMS} = \sqrt{\frac{1}{n}\,\sum_{i=1}^n\,(\mbox{residual error}_i)^2}\vspace{-3pt}\]
It is a measure of the total error of the line that we are using  to fit the data.
The regression line is the line that minimizes this error.  
It is the \textit{best fit} line.
Its RMS can be computed using:\vspace{-3pt}
\[\mbox{RMS}_{\mbox{\footnotesize reg}} = \mbox{sd}_y\,\sqrt{1-r^2}\vspace{-3pt}\]
where $r$ is the correlation and $\mbox{sd}_y$ is the standard deviation of the $y$-values
in our data set.
Notice that:
\begin{itemize}
\item For a fixed value of $r$, RMS increases with $\mbox{sd}_y$.
\item If $r=1$ or $r=-1$, the RMS is zero (since the points fall exactly on a line).
\item If $r=0$, then $\mbox{RMS}=\mbox{sd}_y$.
\end{itemize}

\begin{center}
\begin{pspicture}(0,-0.7)(7,7.7)
\psset{yunit=0.6,dotsize=0.15}
\psdot(1,10)
\psdot(3,9)
\psdot(4,8)
\psdot(4,7)
\psdot(5,1)
\psdot(7,13)
\psline[linewidth=0.02](0,7.8)(7,8.15)(8,8.2)
\rput(3.5,-1.7){$x$}\rput(-1,6.5){$y$}
\psaxes(0,0)(7,13)
%
\psline[linestyle=dotted](7,8.15)(7,13)
\psdot*[linecolor=red](7,8.15)
\rput[l](7,10.575){\ error}
%
\psline[linestyle=dotted](3,7.95)(3,9)
\psdot*[linecolor=red](3,7.95)
\rput[r](3,8.475){error\ }
%
\psline[linestyle=dotted](1,7.85)(1,10)
\psdot*[linecolor=red](1,7.85)
%
\psline[linestyle=dotted](4,8)(4,7)
\psdot*[linecolor=red](4,8)
%
\psline[linestyle=dotted](5,8.05)(5,1)
\psdot*[linecolor=red](5,8.05)
%
\rput(8,5.4){regression line:}
\rput(8,4.5){$y-8\; = \; \frac{1}{20}\,(x - 4)$}
\end{pspicture}
\end{center}

1. Compute the regression line RMS given $\mbox{sd}_y$ and $r$

\hspace{10pt} a) $\mbox{sd}_y=8$ and $r=\frac{\sqrt{3}}{2}$
\vspace{.2in}

\hspace{10pt} b) $\mbox{sd}_y=1$ and $r=\frac{\sqrt{3}}{2}$
\vspace{.2in}

\hspace{10pt} c) $\mbox{sd}_y=1$ and $r=-\frac{\sqrt{3}}{2}$
\vspace{.2in}

\hspace{10pt} d) $\mbox{sd}_y=5$ and $r=0.1$


\vfill
\eject
{\ }

One of the assumptions of the regression model is that for any given $x$, the 
histogram of the possible $y$ values is a normal curve.  
Its mean is the predicted $y$ value (the point on the line).  Its sd is the RMS.  
The figure below shows a normal curve moving along with the regression line.
The $y$-value of each red dot is the mean of the corresponding normal curve.

\begin{center}
\begin{pspicture}(0,-0.7)(7,7.7)
\psset{yunit=0.6,dotsize=0.15}
\psdot(1,10)
\psdot(3,9)
\psdot(4,8)
\psdot(4,7)
\psdot(5,1)
\psdot(7,13)
\psline[linewidth=0.02](0,7.8)(7,8.15)(8,8.2)
\rput(3.5,-1.7){$x$}\rput(-1,6.5){$y$}
\psaxes(0,0)(7,13)
%
\psdot*[linecolor=red](7,8.15)
     \rput{-90}(7,8.15){\psGauss[sigma=0.6]{-2}{2}}
%
\psdot*[linecolor=red](3,7.95)
     \rput{-90}(3,7.95){\psGauss[sigma=0.6]{-2}{2}}
%
\psdot*[linecolor=red](1,7.85)
     \rput{-90}(1,7.85){\psGauss[sigma=0.6]{-2}{2}}
%
\psdot*[linecolor=red](4,8)
     \rput{-90}(4,8){\psGauss[sigma=0.6]{-2}{2}}
%
\psdot*[linecolor=red](5,8.05)
     \rput{-90}(5,8.05){\psGauss[sigma=0.6]{-2}{2}}
%
\end{pspicture}
\end{center}

2. Use the information given to find the mean and sd of the moving normal curve.
If $x=3$,  find a range for 68\% of the $y$-values.

\hspace{10pt} a) $y-8\;=\;2\,(x-4)$, $\mbox{sd}_y=8$, $r=\frac{\sqrt{3}}{2}$.
\vspace{1.5in}

\hspace{10pt} b) $y-2\;=\;5\,(x-1)$, $\mbox{sd}_y=1$, $r=\frac{\sqrt{3}}{2}$.
\vspace{1.5in}

\hspace{10pt} c) $y-8\;=\;-2\,(x-2)$, $\mbox{sd}_y=1$, $r=-\frac{\sqrt{3}}{2}$.


\vfill
\eject

\end{document}

